\chapter{Contact}
\label{chap:contact}

Contact introduces unilateral interaction between a set of slave nodes and a master surface region. Unlike a tie constraint, contact can open and close during the nonlinear solution. FEMaster currently exposes one frictionless formulation: node-to-surface penalty contact for \val{NONLINEARSTATIC} loadcases. The normal law uses a narrow CalculiX-style regularization around zero gap, including the corresponding very small tensile branch on the open side.

\section{Contact definition}

\begin{syntax}
*CONTACT, MASTER=<master-surface>, SLAVE=<slave-surface>, PENALTY=<value>, CLEARANCE=<value>, FLIP=NO|YES
\end{syntax}

\key{MASTER}, \key{SLAVE}, and \key{PENALTY} are required. \key{CLEARANCE} defaults to 0 and \key{FLIP} defaults to \val{NO}. Both names refer to surface sets. The unique nodes contained in the slave surface set become the slave contact nodes; the complete faces contained in the master surface set are the candidate master faces.

\cmd{CONTACT} is defined at model root level. The contact contribution is assembled only during \val{NONLINEARSTATIC}.

\section{Node-to-surface search}

FEMaster follows the topological master search used by CalculiX node-to-face contact. The complete finite-element master faces are first represented by an auxiliary search triangulation. Three-, four-, six-, and eight-node faces are subdivided into 1, 2, 4, and 6 search triangles, respectively. Every search triangle stores the finite-element master face from which it originated. The triangulation is used only for master-face ownership; the final closest-point projection, gap, force, and tangent are evaluated on the complete finite-element face.

Search triangles are connected through common edges. At every master-face node, the normals of the incident master faces are evaluated at that node, accumulated, and normalized. For each search-triangle edge, the normals at its two end nodes are projected perpendicular to the edge and averaged to obtain a representative edge normal. Together with the edge direction, this defines a bordering plane through the edge. The plane normal is oriented outwards from the search triangle. Neighboring triangles sharing an edge therefore use the same geometric separation between their search regions.

For each slave node, the search starts at the master search triangle with the nearest current centroid. The point is tested against the three bordering planes of the current triangle. If it lies outside one plane by more than

\[
    d_\mathrm{search} = 10^{-3}\sqrt{A_s},
\]

where \(A_s\) is the representative slave nodal area, the search crosses that edge to the neighboring triangle and continues there. If no bordering plane is violated, the current triangle owns the slave point. If a step would immediately return to the previously visited triangle, FEMaster follows the CalculiX ``solution in between triangles'' rule and accepts the current triangle. A longer circular walk is treated as an invalid search result.

After a search triangle has been selected, its associated complete master face is used for the ordinary bounded finite-element closest-point projection. The resulting slave-to-master connectivity is retained only when the pair satisfies the positive-gap contact cutoff described below.

\section{Contact topology during Newton iteration}

The search graph determines a discrete contact-element topology: every retained slave node stores the identifier of its current master face. A Newton residual/tangent evaluation may regenerate this topology. The analytic tangent is then formed with exactly those master faces held fixed.

All line-search residual evaluations belonging to that Newton direction reuse the stored topology. They update the continuous closest-point coordinates, gap, normal, force, and master shape functions on the retained face, but they do not run the master search again. Consequently, a slave node cannot switch from one neighboring master face to another while the line search is testing different values of the same Newton correction.

After Newton convergence, the existing nonlinear active-set callback performs one topology-update evaluation at the converged geometry. If the resulting slave-to-master signature differs from the topology used by Newton, the increment is not accepted immediately. Instead, Newton is restarted at the same load factor with the updated contact topology. This is the contact-discontinuity iteration used to make the discrete search state consistent with the converged equilibrium state.

Following CalculiX node-to-face contact, regeneration is allowed through the first eight Newton tangent evaluations of one attempted increment. After the eighth update assembly, the discrete slave-to-master connectivity is frozen for the remaining Newton iterations of that attempt. Continuous geometry and the analytic tangent on each retained master face remain fully current. A cutback rolls the topology back to the last accepted increment and starts the eight-iteration update window again.

The contact state is transactional. A rejected line-search trial restores its parent topology, an accepted line-search trial keeps the same frozen connectivity, and a rejected increment restores the topology from the last accepted physical increment. No distance-based partner-switch hysteresis is used.

\section{Gap and normal}

Let \(\mathbf{x}_s\) denote the current slave-node position, \(\mathbf{x}_m(\xi,\eta)\) the closest point on the selected master face, and \(\mathbf{n}_m(\xi,\eta)\) its current unit normal. FEMaster evaluates the signed normal gap as

\[
    g = \left(\mathbf{x}_s-\mathbf{x}_m\right)\cdot\mathbf{n}_m-c,
\]

where \(c\) is the prescribed \key{CLEARANCE}. Negative gap corresponds to penetration. The contact normal comes from the actual finite-element face interpolation at the closest point. The averaged nodal normals used by the search graph affect only search-region ownership and do not replace the contact normal. If the master normal is opposite to the intended orientation, \key{FLIP=YES} reverses it for the contact pair.

\begin{warningbox}[Surface orientation]
The signed gap depends directly on the master surface orientation. Check the selected element side and its outward normal before changing the penalty parameter. In particular, opposite faces of a solid element must use opposite node winding.
\end{warningbox}

\section{Slave nodal area}

The penalty parameter is interpreted as a pressure-overclosure stiffness. Each slave node therefore receives a representative tributary area assembled from the slave faces that contain it. Linear triangles and quadrilaterals use equal positive nodal shares. For six-node faces the CalculiX weights \(1/999\) at corner nodes and \(332/999\) at midside nodes are used. Eight-node faces use \(1/100\) at corner nodes and \(24/100\) at midside nodes. These positive quadratic-face weights avoid zero or negative representative contact areas.

The nodal characteristic length is

\[
    \ell_s = \sqrt{A_s}.
\]

It defines the search-border tolerance, narrow smoothing, and positive-gap cutoff scales used by the contact formulation. The representative area is held fixed during one contact tangent evaluation, matching the CalculiX contact-spring linearization.

\section{Regularized penalty force}

The linear pressure-overclosure law follows the arctangent regularization used by CalculiX node-to-face contact. FEMaster evaluates

\[
    \mathbf{f}_s
    = \epsilon A_s g
      \left[
          \frac{1}{2}
          + \frac{1}{\pi}
            \arctan\left(-\frac{g}{\delta}\right)
      \right]
      \mathbf{n}_m,
\]

where \(\epsilon\) is the user-specified \key{PENALTY}. The current smoothing width is chosen geometrically as

\[
    \delta = 10^{-6}\ell_s.
\]

For penetration much larger than \(\delta\), the bracket tends to one and the ordinary linear penalty law is recovered. At zero gap the force is exactly zero and the force--gap relation is smooth. For a very small positive gap the regularized law produces the correspondingly small tensile branch also present in the CalculiX formulation.

The contact pair is retained only while

\[
    g \le c_0,
    \qquad
    c_0 = 10^{-3}\ell_s.
\]

Outside this positive-gap cutoff the pair contributes no force and is not retained in the discrete contact-element topology. A node that closes from outside this cutoff during a frozen line-search trial is therefore introduced by the next topology-update evaluation rather than changing the discrete problem inside the line search.

The equal and opposite force is distributed to the master-face nodes using the master shape functions at the closest point:

\[
    \mathbf{f}_{m,a} = -N_a(\xi,\eta)\,\mathbf{f}_s.
\]

Consequently, a closest point on a master edge or corner needs no separate edge-contact formulation: the ordinary finite-element shape functions naturally distribute the master reaction to the nodes active at that location.

\section{Consistent contact tangent}

For the selected master face, FEMaster follows the analytic node-to-face linearization used by CalculiX. Define

\[
    \mathbf{r}=\mathbf{x}_s-\mathbf{x}_m(\xi,\eta),
    \qquad
    \mathbf{a}_\xi=\mathbf{x}_{m,\xi},
    \qquad
    \mathbf{a}_\eta=\mathbf{x}_{m,\eta}.
\]

For an interior closest point the projection satisfies

\[
    \mathbf{r}\cdot\mathbf{a}_\xi=0,
    \qquad
    \mathbf{r}\cdot\mathbf{a}_\eta=0.
\]

Differentiating these equations with respect to each nodal displacement gives a \(2\times2\) system for \(d\xi/du\) and \(d\eta/du\). Its coefficient matrix is

\[
    \begin{bmatrix}
        -\mathbf{a}_\xi\cdot\mathbf{a}_\xi
        +\mathbf{r}\cdot\mathbf{x}_{m,\xi\xi}
        &
        -\mathbf{a}_\xi\cdot\mathbf{a}_\eta
        +\mathbf{r}\cdot\mathbf{x}_{m,\xi\eta}
        \\
        -\mathbf{a}_\xi\cdot\mathbf{a}_\eta
        +\mathbf{r}\cdot\mathbf{x}_{m,\xi\eta}
        &
        -\mathbf{a}_\eta\cdot\mathbf{a}_\eta
        +\mathbf{r}\cdot\mathbf{x}_{m,\eta\eta}
    \end{bmatrix}.
\]

The resulting local-coordinate derivatives are propagated through the relative vector, both master-surface tangents, the normalized master normal and the normal gap. The derivative of the regularized scalar force includes the complete derivative of the arctangent pressure--overclosure law. Finally, the master force linearization also includes

\[
    dN_a=N_{a,\xi}\,d\xi+N_{a,\eta}\,d\eta.
\]

Thus the tangent contains the same geometric terms represented by \(d\xi\), \(d\eta\), the changing surface normal and the changing shape functions in CalculiX \texttt{springstiff\_n2f}. The discrete choice of master face and the representative slave area are not differentiated. In the absence of friction the completed local tangent is symmetrized before sparse assembly, again following the CalculiX node-to-face implementation.

Regression tests independently finite-difference the assembled contact residual against the analytic tangent, verify graph walking across a shared master edge, and verify that a frozen line-search trial keeps the Newton master face while a subsequent topology-update trial detects the corresponding partner switch.

\section{Current limitations}

The current contact implementation is intentionally narrow:

\begin{itemize}
\item frictionless normal contact only,
\item slave nodes obtained from a slave surface set,
\item master finite-element faces obtained from a master surface set,
\item CalculiX-style topological search triangulation without a spatial acceleration structure,
\item regularized penalty enforcement only,
\item contact contributions only in \val{NONLINEARSTATIC},
\item no surface-to-surface or mortar formulation,
\item no augmented-Lagrange multipliers,
\item no BVH or other spatial acceleration structure,
\item no distance-based partner-switch hysteresis,
\item no tangential friction law.
\end{itemize}

\section{Minimal example}

\begin{exampledeck}
*SURFACE, NAME=MASTER_FACE
1, 1, S2
*SURFACE, NAME=SLAVE_FACE
2, 2, S1

*CONTACT, MASTER=MASTER_FACE, SLAVE=SLAVE_FACE, PENALTY=1.0e5

*LOADCASE, TYPE=NONLINEARSTATIC, NAME=CONTACT_STEP
*SUPPORTS
BC
*SOLVER, DEVICE=CPU, METHOD=DIRECT
*CONSTRAINTMETHOD, TYPE=NULLSPACE
*NONLINEAR, CONTROL=LOAD, INITIAL_INCREMENT=0.1, MAXIMUM_INCREMENT=0.1, MAX_INCREMENTS=10
*END
\end{exampledeck}

A complete two-block regression model is included in \file{examples/10_contact/103_contact_node_surface_parallel_blocks_c3d8.inp}.
